\documentclass[11pt,a4paper]{article}
\usepackage[T1]{fontenc}
\usepackage[utf8]{inputenc}
\usepackage{lmodern,amsmath,amssymb,amsthm,geometry,microtype,booktabs,enumitem}
\usepackage[hidelinks]{hyperref}
\geometry{margin=25mm}
\setlength{\parskip}{4pt}
\setlength{\parindent}{0pt}
\setlist{nosep,leftmargin=*}
\newcommand{\ceil}[1]{\left\lceil #1\right\rceil}
\newcommand{\floor}[1]{\left\lfloor #1\right\rfloor}
\newcommand{\pt}{\operatorname{pt}}
\newcommand{\thone}{\operatorname{th}_1}
\newtheorem*{theoremA}{Theorem A (fixed-set worst time)}
\newtheorem*{theoremB}{Theorem B (fixed-budget optimum)}
\newtheorem*{theoremC}{Theorem C (additive throttling)}
\newtheorem*{corB1}{Corollary B1 (deadline certificate)}
\newtheorem*{corB2}{Corollary B2 (minimum seeds at a deadline)}
\newtheorem*{lemmaO1}{Observation O1 (initial blue--blue edge removal)}
\title{Exact Parallel Propagation and Throttling\\with One Leak on a Cycle}
\author{Ryutaro Yonezu\\Yonezu Research Lab\\
\small\href{mailto:yonezaemon1+research@gmail.com}{yonezaemon1+research@gmail.com}}
\date{26 October 2026\quad\textbar\quad Version 1.0\\
\small Research Note No.41\quad\textbar\quad Unrefereed preprint}
\begin{document}
\maketitle
\begin{abstract}
Consider standard zero forcing on a simple undirected cycle, with simultaneous
rounds and one fixed adversarial vertex that cannot force but may become blue.
For an arbitrary initial set, retain its maximal blue runs of length at least
two. The worst completion time is the largest number of initially white
vertices in an arc between consecutive retained runs; isolated initial blue
vertices do not split such an arc. We give an elementary proof using the known
one-leaky robustness of path endpoints, a time-preserving removal of initial
blue--blue edges, and an explicit lower-bound witness. Balanced layouts then
give the exact seed-budget frontier, deadline minimum, and additive throttling
value $\ceil{\sqrt{8n}}-2$. The established forcing model and propagation
mechanism are not new. This note makes an exact time formula explicit;
historical novelty remains unestablished.
\end{abstract}

\section{Scope and relation to established results}

Leaky forcing, ordinary propagation time, and throttling are established
notions. Alameda, Dillman and Kenter prove that the endpoints of a path form a
one-leaky forcing set \cite[Proposition 3.1, p.312]{ADK2024}. The same fact
follows by specializing the tree theorem of Alameda, Kritschgau, Warnberg and
Young \cite[Theorem 4.4]{AKWY2022} to one leak. The usual cycle robustness
argument uses two opposing fronts \cite[Proposition 3.4, p.313]{ADK2024}.
We use these facts rather than claim a new fault-bypass mechanism.

The upper time count below uses the familiar principle that a successful
process must make progress at every step. Hogben et al. use this principle
for an ordinary propagation-time bound \cite[Remark 1.8]{Hogben2012}.
Here its application to an arbitrary initial set and a fixed leak is justified
explicitly. Butler and Young develop the forcing-chain capacity and path
construction behind standard throttling \cite[Section 2, pp.66--68]{BY2013}.
The arithmetic optimization in this note is a consequence of the fixed-set
formula, not an independent optimization principle.

Carlson et al. determine ordinary cycle throttling and study PSD and weighted
PSD throttling \cite[Sections 6--7]{Carlson2019}. Their color-change rules
must be distinguished from the fixed vertex-leak rule used here. In
particular, standard forcing does not let an isolated initial blue vertex
start a front while both of its neighbors remain white.

The intended scope is an explicit exact time formula and an elementary
reproof using known path facts plus the supplementary observations proved
below. No priority, new forcing model, human peer review, or complete
literature search is claimed. An elementary derivation does not establish
that the same time formula was previously published. This is a mathematical
graph model, without a claim about measured physical transmission time.

\section{Fixed model and notation}

Let $C_n$ be the simple undirected cycle on $\mathbb Z/n\mathbb Z$, $n\geq3$,
with edges $\{i,i+1\}$. First choose an initial blue set $B$. An adversary then
chooses one fixed vertex $x$ from \emph{all} vertices, including initially
white ones. Vertex $x$ cannot force, but can itself become blue. No vertex or
edge is deleted in the model.

Time zero is the initial coloring. In each round, every blue vertex other
than $x$ with exactly one white neighbor forces that neighbor blue. All
decisions use the coloring at the \emph{start} of the round and are applied
simultaneously. A newly blue vertex can force only in a later round.
There is no domination step, hopping, PSD component rule, or moving leak.

Let $p_G(B;x)$ be the first round at which all vertices of $G$ are blue, or
$\infty$ if the process stalls, and set
\begin{align*}
 T_G(B)&=\max_{x\in V(G)}p_G(B;x),\\
 \tau(n,k)&=\min_{|B|=k}T_{C_n}(B),\\
 m(n,t)&=\min\{|B|:T_{C_n}(B)\leq t\},\\
 \thone(C_n)&=\min_B\bigl(|B|+T_{C_n}(B)\bigr).
\end{align*}
Here $t$ is a nonnegative integer, and infinite-time sets do not improve the
last minimum. We write $T(B)=T_{C_n}(B)$ when clear. For $k=n$ the time is
zero; for $k=0,1$ it is infinite. The notation $\thone$ is local to this note.

For a fixed graph and leak, enlarging the initial blue set preserves the
round-by-round inclusion of blue sets. Indeed, a force $u\to v$ valid for
a smaller blue set is either already satisfied in the larger set or remains
valid because $u$ and all its other neighbors are blue there. This also shows
that applying all available forces is at least as fast as any legal schedule.

\section{The fixed-set formula}

For a proper initial set $B$, let $R_1,\ldots,R_r$ be its maximal initial
blue runs of length at least two, in cyclic order. For $r>0$, let $I_i$ be
the arc strictly between $R_i$ and $R_{i+1}$, with cyclic indices. Let $a_i$
be the last vertex of $R_i$ before the arc and $b_i$ the first vertex of
$R_{i+1}$ after it. These endpoints are distinct even when $r=1$.

The arc begins and ends with white vertices. Any initial blue vertices
inside it are isolated. Let $W_i$ count the \emph{initially white} vertices
in $I_i$. Thus $W_i\geq1$ and
\begin{equation}
 \sum_i W_i=n-|B|. \label{eq:whitecount}
\end{equation}
Its length may exceed $W_i$.

\begin{lemmaO1}
Let $F$ be any collection of edges of a graph $G$ whose endpoints are both
initially blue, and put $H=G-F$. Keep every vertex, the initial set, and the
fixed leak unchanged. Then the blue set at every round is identical in $G$
and $H$. In particular $p_H(B;x)=p_G(B;x)$ for every $x$.
\end{lemmaO1}
\begin{proof}
For any current blue set $D\supseteq B$ and any vertex $v$,
\[
 N_H(v)\setminus D=N_G(v)\setminus D.
\]
Every removed neighbor is in $B$ and remains blue. No removed edge can be
used to force a white target. Hence the possible forces, including the
restriction at $x$, are the same. Induction from time zero proves the claim.
\end{proof}

\begin{theoremA}
If $B=V(C_n)$ then $T(B)=0$. If $B$ is proper with no adjacent initial blue
vertices then $T(B)=\infty$. In all other cases,
\begin{equation}
 T(B)=\max_i W_i. \tag{A}\label{eq:A}
\end{equation}
Placing the leak at $b_j$ for any arc with maximum $W_j$ attains this time.
\end{theoremA}
\begin{proof}
The full-set case is immediate. With no adjacent initial blue vertices,
every blue vertex has two white neighbors, so no force starts, even without
a leak.

Otherwise apply O1 to all initial blue--blue edges:
\[
 H=C_n-E(C_n[B]).
\]
The components containing white vertices are disjoint paths
$P_i=a_i,I_i,b_i$, with both endpoints initially blue and only isolated
initial blue vertices inside. Other components are isolated blue vertices.
This description also covers a run of length two and the case of one long
run. It preserves the original $W_i$. The edge removal is an analytical
transformation under O1, not a change to the failure model.

For the upper bound, the two endpoints of each $P_i$ form a one-leaky forcing
set by the known path result \cite[Proposition 3.1]{ADK2024}, or the tree
theorem \cite[Theorem 4.4]{AKWY2022}. Adding the internal initial blue
vertices preserves success by monotonicity. If the global leak is outside
the component, its process has no leak and succeeds as well.

A successful fixed-leak process cannot have a round with no color change
before completion. Such a round would leave the same coloring and the same
prohibition, causing permanent stalling. Each incomplete round therefore
colors at least one of the $W_i$ white vertices. Thus $P_i$ finishes within
$W_i$ rounds. The components run in parallel, and O1 gives
$T(B)\leq\max_i W_i$.

For the lower bound, choose $P_j$ with maximum $W_j$ and leak its endpoint
$b_j$. List its initially white vertices from $a_j$ to $b_j$ as
$w_1,\ldots,w_{W_j}$. We show that at time $q$, $0\leq q\leq W_j$,
precisely $w_1,\ldots,w_q$ have become blue. This is true at time zero.
If $q<W_j$, the vertex immediately before $w_{q+1}$ is a nonleaky blue
vertex: it is $a_j$, $w_q$, or the initial blue singleton between $w_q$ and
$w_{q+1}$. Its backward neighbor, if present, is blue. It can therefore
force $w_{q+1}$ in the next round.

Every initial blue singleton further along still has two white neighbors
at the start of that round. In particular, the singleton just beyond
$w_{q+1}$ cannot simultaneously force the following white vertex: $w_{q+1}$
is still white when decisions are made. Newly colored vertices cannot force
within the same round, and entry from $b_j$ is prohibited. Thus no other
white target is reached. This proves the induction. The component needs
exactly $W_j$ rounds, and O1 gives $T(B)\geq W_j$.
\end{proof}

O1 and the lower-bound count are supplementary elementary arguments proved
here. The cited reachability theorem does not itself state (A).

\section{Seed budgets and deadlines}

\begin{theoremB}
For $n\geq3$ and $2\leq k<n$,
\begin{equation}
 \tau(n,k)=\ceil{\frac{n-k}{\floor{k/2}}}. \tag{B}\label{eq:B}
\end{equation}
The remaining values are $\tau(n,n)=0$ and
$\tau(n,0)=\tau(n,1)=\infty$.
\end{theoremB}
\begin{proof}
For a finite-time proper set of size $k$, every long blue run contains at
least two seeds, so $r\leq\floor{k/2}$. Equations (A) and
\eqref{eq:whitecount} give
\[
 n-k=\sum_iW_i\leq rT(B)\leq\floor{k/2}T(B),
\]
which proves the lower bound, including sets with isolated seeds.

For the upper bound, put $W=n-k>0$ and
$r=\min\{\floor{k/2},W\}$. Form $r$ blue runs of at least two seeds each,
put all $k-2r$ extra seeds in the first run, and separate them by $r$
nonempty white runs whose lengths sum to $W$ and differ by at most one.
There are no isolated blue seeds in this layout. If $W\geq\floor{k/2}$,
(A) gives the desired ceiling. If $W<\floor{k/2}$, each white run has one
vertex and the time is one, also the value in (B).
\end{proof}

\begin{corB1}
For an integer $t\geq0$, a proper $B$ completes by $t$ against every fixed
leak if and only if it has a long initial blue run and each intervening arc
contains at most $t$ initially white vertices. The full set completes at zero.
\end{corB1}
\begin{proof}
This is exactly (A), including its stalled case.
\end{proof}

\begin{corB2}
For $n\geq3$ and an integer $t\geq0$,
\begin{equation}
 m(n,t)=\min\left\{2\ceil{\frac{n}{t+2}},\;
  2\ceil{\frac{n-1}{t+2}}+1\right\}. \tag{C}\label{eq:C}
\end{equation}
In particular $m(n,0)=n$, and $m(n,t)=2$ if and only if $t\geq n-2$.
\end{corB2}
\begin{proof}
By (B), a budget $k\geq2$ is feasible precisely when
$n\leq k+t\floor{k/2}$, including $k=n$. For even $k=2r$, the least
candidate is $2\ceil{n/(t+2)}$; for odd $k=2r+1$, it is
$2\ceil{(n-1)/(t+2)}+1$. A candidate above $n$ cannot win: for $t=0$
one parity gives exactly $n$, and for $t\geq1$ the even candidate is at
most $n$. The odd candidate is at least three and does not introduce $k=1$.
Finally $k=0,1$ are infeasible and $\tau(n,2)=n-2$.
\end{proof}

\section{Additive throttling}

\begin{theoremC}
For all $n\geq3$,
\begin{equation}
 \thone(C_n)=\ceil{\sqrt{8n}}-2. \tag{D}\label{eq:D}
\end{equation}
\end{theoremC}
\begin{proof}
Let $c$ be an integer cost budget, $k$ the seed count, and $t$ the time.
Feasibility with $k+t\leq c$ implies
\[
 n\leq k+(c-k)\floor{k/2}.
\]
For $k=2r$ the right side is $r(c+2-2r)$; for $k=2r+1$ it is
$r(c+1-2r)+1$. The even expression exceeds the odd one by $r-1\geq0$.
The cases $k=0,1$ are infeasible. For $c\geq3$,
\begin{equation}
 \max_{1\leq r\leq\floor{c/2}} r(c+2-2r)
       =\floor{\frac{(c+2)^2}{8}}. \tag{E}\label{eq:E}
\end{equation}
Indeed, setting $h=c+2$ gives
\[
 r(h-2r)=\frac{h^2-(h-4r)^2}{8}.
\]
Choose an integer $r$ nearest to $h/4$. The squared loss is $0$, $1$, or
$4$ when $h\equiv0$, $h$ is odd, or $h\equiv2\pmod4$, respectively,
giving exactly the displayed floor. Such a maximizing $r$ lies in the
required range for $c\geq3$.

It follows that $c\geq\ceil{\sqrt{8n}}-2$. Conversely set
$c=\ceil{\sqrt{8n}}-2$. For $n\geq3$, $3\leq c\leq n$ because
$\sqrt{8n}\leq n+2$. Choose a maximizing $r$ in (E) and set
$k=2r$, $t=c-2r$. Then $2\leq k\leq c\leq n$ and $t\geq0$.
The inequality $n\leq r(c+2-2r)$ is exactly $n-k\leq rt$.
The balanced layout in Theorem B therefore attains time at most $t$ and
total cost at most $c$.
\end{proof}

This is elementary capacity arithmetic after (A) and (B). It is not offered
as a new root law independent of the established throttling framework.

\section{Examples and model checks}

On $C_5$, take $B=\{0,1,3\}$. Every contiguous white run has length one,
but there is a single arc between long blue runs, containing two white
vertices. With leak $x=0$, the newly blue sets are $\{2\}$ and then
$\{4\}$, so $T(B)=2$. Counting only the longest contiguous white run
would give a false answer.

On $C_{100}$ with $k=14$, use seven blue pairs and white gaps of lengths
$13,13,12,12,12,12,12$. The worst time is $13$; (B) rules out time $12$
for every $14$-seed layout. The sum $14+13=27$ attains (D).

On $C_6$ with $B=\{0,1\}$, ordinary forcing without a leak finishes in
two rounds, whereas the worst one-leak time is four. Removing one initial
seed instead leaves a single blue vertex and stalls. Initial-seed loss is
therefore not interchangeable with this vertex-leak model.

If two or more arbitrary vertex leaks are allowed and any vertex is omitted
from $B$, the adversary can leak both its neighbors. The omitted vertex can
never be colored. Thus only the full set is robust. This is the established
low-degree obstruction \cite[Lemma 2.4]{ADK2024}, not an additional result.

\section{Verification, provenance, and limits}

The accompanying Python verifier directly simulates the literal rule for
every labeled initial set and every fixed leak. It compares the results
with (A), checks the adversarial endpoint, and exhaustively minimizes over
all sets to check (B), (C), and (D). Its recorded author-side run is:
\begin{center}
\begin{tabular}{lr}
\toprule
Quantity & Recorded value\\
\midrule
Range & $3\leq n\leq16$\\
Initial sets & $131{,}064$\\
Set/leak scenarios & $1{,}966{,}072$\\
Fixed-budget entries, including infeasible sizes & $147$\\
Deadline entries & $147$\\
Mismatches & $0$\\
\bottomrule
\end{tabular}
\end{center}
Construction checks and model-discrimination examples are additional and are
not included in the exhaustive scenario count. The recorded environment was
Python 3.12.14 on Linux, approximately 15 seconds. Verification uses the
Python standard library, with no network, paid API, GPU, or cloud computation
required. Windows execution has not been tested by the author-side record.
Finite tests are falsification checks, not proofs for unbounded $n$.

User-supplied self-audit reports find no model, proof, or computational
defect and report a separately written C check through $n=22$. Its C source
and raw logs were not supplied to this package, so that run is retained as
a reported result, not authenticated author-side evidence. The corrected
reported cumulative counts are $8{,}388{,}600$ initial sets and
$176{,}160{,}760$ scenarios for $3\leq n\leq22$; the earlier approximate
totals referred to $n=21,22$ only. Repeated reports from the same audit
process are not independent third-party confirmations. No human peer
review is claimed.

The evidence archive preserves candidate v0.4 and its source-location
checks. Selected source pages were checked by the author-side workflow;
overall literature comparison remains partial. The unresolved question is
historical originality of the exact fixed-set time statement, including
equivalent formulations and direct consequences of earlier work. The
publication makes the explicit elementary formula available for inspection
without asserting priority or a new propagation principle.

\paragraph{Availability and rights.}
The paper source, verification materials, and audit snapshot are available
in the accompanying evidence archive and at
\url{https://github.com/yonezaemon1-hub/no41-one-leak-cycle-propagation}.
Prepared with AI assistance. Copyright \copyright\ 2026 Ryutaro Yonezu.
All rights reserved. Research inquiries:
\href{mailto:yonezaemon1+research@gmail.com}{yonezaemon1+research@gmail.com}.

\begin{thebibliography}{9}
\bibitem{ADK2024}
J. S. Alameda, S. Dillman and F. Kenter,
\emph{Leaky forcing: a new variation of zero forcing},
Australasian Journal of Combinatorics \textbf{88}(3) (2024), 308--326.
\url{https://ajc.maths.uq.edu.au/pdf/88/ajc_v88_p308.pdf}.

\bibitem{AKWY2022}
J. S. Alameda, J. Kritschgau, N. Warnberg and M. Young,
\emph{On leaky forcing and resilience},
Discrete Applied Mathematics \textbf{306} (2022), 32--45.
\href{https://doi.org/10.1016/j.dam.2021.09.005}{doi:10.1016/j.dam.2021.09.005}.
Inspected author version: \url{https://arxiv.org/abs/2008.06552v1},
Theorem 4.4, PDF p.13.

\bibitem{Hogben2012}
L. Hogben, M. Huynh, N. Kingsley, S. Meyer, S. Walker and M. Young,
\emph{Propagation time for zero forcing on a graph},
Discrete Applied Mathematics \textbf{160} (2012), 1994--2005.
\href{https://doi.org/10.1016/j.dam.2012.04.003}{doi:10.1016/j.dam.2012.04.003}.
Inspected author PDF: \url{https://arxiv.org/abs/1410.4191v1},
Remark 1.8, PDF p.3; the delivered PDF has a 2018 body date.

\bibitem{BY2013}
S. Butler and M. Young,
\emph{Throttling zero forcing propagation speed on graphs},
Australasian Journal of Combinatorics \textbf{57} (2013), 65--71.
\url{https://ajc.maths.uq.edu.au/pdf/57/ajc_v57_p065.pdf}.

\bibitem{Carlson2019}
J. Carlson, L. Hogben, J. Kritschgau, K. Lorenzen, M. S. Ross,
S. Selken and V. Valle Martinez,
\emph{Throttling positive semidefinite zero forcing propagation time on graphs},
Discrete Applied Mathematics \textbf{254} (2019), 33--46.
\href{https://doi.org/10.1016/j.dam.2018.06.017}{doi:10.1016/j.dam.2018.06.017}.
Inspected author version: \url{https://arxiv.org/abs/1807.09404v1},
Theorem 7.1 and Lemma 7.2, PDF p.19.
\end{thebibliography}
\end{document}
